\chapter{When Things Go Wrong — Lessons from Real Cases}
\label{chap:When Things Go Wrong — Lessons from Real Cases}

%---------------------------- SECTION ----------------------------
\section{The Value of Failure}
\paragraph{There is a particular kind of knowledge that cannot be acquired from frameworks, methodologies, or theoretical analysis alone. It is the knowledge that comes from studying what actually happened — from examining, with honesty and rigour, the cases where risk assessment failed, where warnings were ignored, where controls collapsed, and where the consequences fell on real people, real organisations, and real communities.}
\paragraph{This knowledge is uncomfortable. Real cases involve real failures by real people — often intelligent, experienced, well-intentioned people operating within organisations that had risk management frameworks, compliance functions, and governance structures that were, on paper, adequate. The discomfort of examining those failures honestly — without the consolation of attributing them to obvious incompetence or deliberate wrongdoing — is precisely what makes it valuable.}
\paragraph{The cases examined in this chapter have been chosen not for their notoriety alone but for what they reveal about the recurring patterns of risk assessment failure — the dynamics that appear, in different forms and different sectors, in case after case. They are chosen because their lessons are transferable — because the compliance professional who understands them is better equipped to identify the same dynamics in their own organisation before they produce the same outcomes.}
\paragraph{Each case is examined through the analytical lens developed throughout this book — asking not just what went wrong but why, and what a different approach to risk assessment, governance, or culture might have produced a different result.}

%---------------------------- SECTION ----------------------------
\section{Case 1 — The 2008 Financial Crisis: The Failure of Quantitative Confidence}

\subsection{What Happened}
\paragraph{The 2008 global financial crisis remains the most significant and most studied failure of institutional risk assessment in living memory. Its causes were multiple, complex, and deeply interconnected — but at its heart was a catastrophic failure of risk assessment across financial institutions, rating agencies, regulators, and policymakers that allowed an enormous accumulation of systemic risk to go unrecognised until it was too late.}
\paragraph{The proximate mechanism was the collapse of the US subprime mortgage market and the complex financial instruments — collateralised debt obligations (CDOs), mortgage-backed securities, and their derivatives — through which the risk of that market had been distributed, concentrated, and obscured across the global financial system. When US house prices began to fall and subprime mortgage defaults accelerated, the value of these instruments collapsed — triggering losses that spread rapidly through institutions worldwide, freezing credit markets, and producing the deepest global recession since the 1930s.}

\subsection{The Risk Assessment Failures}
\paragraph{\textbf{The quantitative model failure}. The risk models on which financial institutions and rating agencies relied — Value at Risk models, credit rating models, CDO valuation models — shared a common and fatal flaw: they were built on historical data that did not include a scenario in which US house prices fell simultaneously across all geographic markets. The models assumed that regional housing markets were not perfectly correlated — and that assumption had been broadly correct historically. But the combination of loose lending standards, securitisation that removed the originator's incentive to maintain credit quality, and the systemic concentration of risk in instruments whose complexity obscured their true exposure created conditions that had no historical precedent.}
\paragraph{The models produced confident outputs — precise risk ratings, narrow confidence intervals, reassuring capital adequacy calculations — from assumptions that were fundamentally wrong for the conditions that actually developed. This is model uncertainty in its most consequential form: not imprecise parameters but structurally misspecified models producing false precision in the domain where precision mattered most.}
\paragraph{\textbf{The complexity and opacity failure}. The instruments through which mortgage risk had been distributed and repackaged — CDO squared structures, synthetic CDOs, credit default swaps on CDO tranches — were so complex that the ultimate risk exposures embedded within them were opaque even to sophisticated market participants. Risk assessment frameworks that were adequate for straightforward instruments were wholly inadequate for instruments whose risk characteristics could only be understood through modelling that itself depended on the flawed assumptions described above.}
\paragraph{Complexity is a risk multiplier — it amplifies the consequences of model errors by making those errors harder to detect and harder to correct. Organisations that accepted complexity as inevitable — rather than treating it as itself a risk factor that warranted additional scrutiny and caution — were systematically blind to the scale of the risk they were carrying.}
\paragraph{\textbf{The incentive failure}. The incentive structures throughout the mortgage origination and securitisation chain systematically undermined the quality of risk assessment. Mortgage originators who sold loans into securitisation vehicles had no long-term stake in the credit quality of those loans. Rating agencies were paid by the issuers of the instruments they rated — creating conflicts of interest that are remarkable in retrospect for their transparency. And traders and structurers who created and distributed complex instruments were rewarded on the basis of short-term revenues rather than long-term risk-adjusted returns.}
\paragraph{As discussed in Chapter~\ref{chap:Building a Risk-Aware Culture}, incentive misalignment is one of the most powerful cultural determinants of risk management quality. An incentive structure that rewards risk-taking without adequate accountability for risk outcomes will consistently produce under-assessment of risk — regardless of the sophistication of the formal risk assessment framework.}
\paragraph{\textbf{The governance failure}. Board-level engagement with the complexity and scale of the risks being taken was, in many institutions, wholly inadequate. Boards that approved business strategies and risk appetites without genuine understanding of the instruments through which those strategies were being executed — and without the capability or the willingness to challenge the quantitative assurances they were given — failed in their most fundamental governance responsibility.}

\subsection{The Lessons}
\paragraph{\textbf{Quantitative confidence is not the same as analytical reliability}. The precision of a risk model's outputs is determined by the quality of its assumptions — and assumptions that are not stress-tested against genuinely adverse scenarios, that are not challenged by independent scrutiny, and that are not subjected to the question \textit{"what would have to be true for this model to be wrong?"} are a source of false confidence rather than genuine understanding.}
\paragraph{\textbf{Complexity is a risk factor, not just a feature}. Instruments, processes, and systems whose risk characteristics cannot be clearly understood and articulated by those who own and manage them should be treated as carrying additional risk — not as neutral features of a sophisticated operation. The inability to explain clearly how a risk instrument works is itself a risk indicator.}
\paragraph{\textbf{Incentive misalignment undermines risk assessment regardless of framework quality}. An organisation whose incentive structures reward risk-taking without adequate accountability for outcomes will systematically produce optimistic risk assessments — because the people producing those assessments have a personal interest in their optimism. No risk assessment framework can fully compensate for pervasive incentive misalignment.}
\paragraph{\textbf{Historical data is not a reliable guide to unprecedented conditions}. Risk models built on historical data will fail precisely in the conditions — unprecedented, extreme, correlated — that matter most. Stress testing must go beyond historical scenarios to explore genuinely novel combinations of adverse conditions.}

%---------------------------- SECTION ----------------------------
\section{Case 2 — The Challenger Space Shuttle Disaster: Normalisation of Deviance}

\subsection{What Happened}
\paragraph{On 28 January 1986, the Space Shuttle Challenger broke apart 73 seconds into its flight, killing all seven crew members. The immediate technical cause was the failure of O-ring seals in a solid rocket booster joint — allowing hot gases to escape and ignite the external fuel tank. The O-ring failure was caused by the unusually cold temperatures at the launch site that morning.}
\paragraph{What made the Challenger disaster one of the most studied cases in risk management history was not the technical failure itself — it was what the subsequent investigation revealed about how the decision to launch had been made, and why the concerns of engineers who understood the O-ring risk had been overridden.}

\subsection{The Risk Assessment Failures}
\paragraph{\textbf{Normalisation of deviance}. NASA engineer Diane Vaughan's investigation of the disaster — described in her landmark study The Challenger Launch Decision — revealed the concept that became one of the most important contributions to risk management thinking: normalisation of deviance. Over the course of the shuttle programme, engineers had repeatedly observed O-ring erosion at temperatures lower than the design specification. But because each flight with O-ring erosion had not resulted in catastrophic failure, the erosion was progressively reframed as an acceptable — and then a normal — feature of shuttle operations rather than a warning signal requiring action.}
\paragraph{This is the accumulation problem in its most consequential form. Each individual instance of O-ring erosion was assessed in isolation — and in isolation, because no catastrophic failure had occurred, it appeared manageable. The pattern — the trend of increasing erosion at lower temperatures — was never assessed as a whole. The deviation from the design standard had been normalised through repetition, and the risk had been systematically under-assessed as a result.}
\paragraph{\textbf{The management override of technical judgement}. The night before the Challenger launch, engineers at Morton Thiokol — the contractor responsible for the solid rocket boosters — made an explicit recommendation against launch at the forecast temperatures. Their recommendation was overridden by managers who were under pressure from NASA to maintain the launch schedule. The engineers' data was challenged, their reasoning was questioned, and the burden of proof was, in the words of one NASA manager, reversed — from \textit{"prove it is safe to fly"} to \textit{"prove it is unsafe to fly"}.}
\paragraph{This reversal of the burden of proof under schedule and commercial pressure is a recurring dynamic in serious incidents across many sectors. When the default changes from \textit{"we launch unless safety is demonstrated"} to \textit{"we delay only if unsafety is demonstrated"}, the threshold for action on safety concerns is raised precisely when it should be at its lowest.}
\paragraph{\textbf{The communication failure}. The engineers' concerns about O-ring performance at low temperatures were communicated within the contractor organisation but were not adequately conveyed to the NASA decision-makers who had the authority to delay the launch. The information existed — but the governance and communication structures through which it should have reached the decision-makers were inadequate, and the culture did not support the kind of persistent, uncomfortable challenge that the situation required.}

\subsection{The Lessons}
\paragraph{\textbf{Normalisation of deviance is a slow and invisible process}. It does not announce itself. It occurs incrementally, through the repeated experience of deviation without immediate adverse consequence — and it can proceed to an advanced stage before anyone has consciously recognised it. Detecting normalisation of deviance requires systematic monitoring for deviations from standards — not just from outcomes — and a culture in which the absence of harm is not treated as evidence that deviation is safe.}
\paragraph{\textbf{The burden of proof matters enormously}. Where the default is action unless risk is demonstrated, organisations are resilient to risk under-assessment — if the risk assessment is wrong, the default protects. Where the default is inaction unless risk is demonstrated, organisations are fragile — if the risk assessment is wrong, or if the communication of concerns is imperfect, the default fails to protect. For significant safety and compliance risks, the burden of proof should be structured to favour caution.}
\paragraph{\textbf{Technical expertise must have genuine voice in decisions that depend on it}. A decision-making structure that systematically overrides technical expertise under schedule or commercial pressure is a risk governance failure — regardless of how well-documented the override process is. Governance structures need to genuinely protect the ability of those with relevant expertise to be heard, to persist in their concerns, and to escalate when they are not adequately addressed.}

%---------------------------- SECTION ----------------------------
\section{Case 3 — The Grenfell Tower Fire: Systemic Failures in Building Safety Risk Assessment}

\subsection{What Happened}
\paragraph{On 14 June 2017, a fire broke out in Grenfell Tower, a 24-storey residential block in North Kensington, London. The fire spread with catastrophic speed, enabled by the aluminium composite material (ACM) cladding that had been installed during a recent refurbishment. Seventy-two people died — the greatest loss of life in a structural fire in the UK since the Second World War.}
\paragraph{The Grenfell Tower Inquiry — whose findings have been published in phases since 2019 — has revealed a systemic failure of building safety risk assessment that operated across multiple organisations, multiple regulatory frameworks, and multiple decades.}

\subsection{The Risk Assessment Failures}
\paragraph{\textbf{The regulatory gap}. The Building Regulations applicable to high-rise residential buildings required that the external walls of such buildings adequately resist the spread of fire. But the regulatory framework contained ambiguities about the standard applicable to cladding systems — and those ambiguities were exploited, interpreted, and in some cases deliberately manipulated to enable the use of materials that did not meet the spirit of the regulatory requirement.}
\paragraph{The gap between the regulatory standard as written and the regulatory standard as genuinely intended is a recurring theme in building safety failures — and it reflects a fundamental tension between a regulatory framework that sets principles and a construction industry that seeks the most commercially convenient interpretation of those principles. Risk assessment that takes the most permissive interpretation of a regulatory requirement — rather than genuinely asking whether the spirit of the requirement is met — is risk assessment in the service of commercial optimisation rather than genuine safety management.}
\paragraph{\textbf{The testing and certification failure}. The ACM cladding used on Grenfell Tower had been tested and certificated — but the tests had been conducted under conditions that did not represent the way the cladding was actually installed and used. The certification process had been exploited — testing configurations that would pass rather than testing configurations that reflected real-world installation. Risk assessments that relied on the existence of test certificates without scrutinising the basis of those certificates were, in effect, relying on evidence that did not support the conclusions drawn from it.}
\paragraph{\textbf{The fragmentation of accountability}. The refurbishment of Grenfell Tower involved multiple contractors, sub-contractors, consultants, and public bodies — each responsible for a specific element of the work, none responsible for the overall safety of the building as a system. This fragmentation of accountability created gaps — responsibilities that fell between the different parties, risks that were not assessed because each party assumed another was responsible for them.}
\paragraph{The fragmentation of accountability problem is a systemic feature of complex construction and infrastructure projects — and it requires specific governance interventions to address. The CDM Regulations' Principal Designer and Principal Contractor roles are precisely designed to address it — but their effectiveness depends on those roles being adequately resourced, genuinely empowered, and actively exercised.}
\paragraph{\textbf{The resident voice ignored}. Residents of Grenfell Tower had raised concerns about fire safety for years before the disaster — through formal complaints, through the Grenfell Action Group blog, and through direct communication with the building manager. Those concerns were not adequately investigated, not adequately responded to, and not incorporated into the building's fire risk assessment. The people with the most direct and most personal stake in the building's safety were the people whose knowledge was most systematically discounted.}

\subsection{The Lessons}
\paragraph{\textbf{Regulatory compliance is not the same as safety}. An organisation that manages its risk assessment to the minimum regulatory standard — that asks "does this comply?" rather than "is this safe?" — is systematically exposed to the gap between what regulation requires and what genuine safety demands. The most dangerous regulatory environments are those where the gap between regulatory compliance and genuine safety is widest — and those are precisely the environments in which commercially motivated interpretation of regulatory requirements is most prevalent.}
\paragraph{\textbf{Test certificates and third-party assurances require scrutiny, not acceptance}. The existence of a test certificate, a compliance declaration, or a third-party assurance is not the same as evidence of safety. Risk assessments that accept such documents at face value without scrutinising the basis on which they were issued are accepting apparent assurance rather than genuine assurance.}
\paragraph{\textbf{Fragmented accountability creates systemic risk gaps}. Complex projects and operations involving multiple parties require explicit governance mechanisms to ensure that overall safety accountability is clearly assigned — and that the gaps between individual party responsibilities are specifically addressed rather than assumed to be covered.}
\paragraph{\textbf{Those closest to the risk have the most valuable risk intelligence}. The consistent failure to engage with, and act on, the concerns of residents who were living with the building's fire safety risks is both a governance failure and a moral one. In any context, the people most directly exposed to a risk have risk intelligence that the formal assessment process should actively seek — not discount.}

%---------------------------- SECTION ----------------------------
\section{Case 4 — The PPI Mis-Selling Scandal: Culture and Incentives Driving Conduct Failure}

\subsection{What Happened}
\paragraph{Payment Protection Insurance (PPI) was sold by UK banks, building societies, and other lenders alongside credit products — mortgages, personal loans, credit cards — from the 1990s through to its effective end following regulatory intervention in the 2010s. PPI was supposed to protect borrowers who were unable to meet their repayments due to illness, accident, or unemployment.}
\paragraph{In practice, PPI was the subject of the largest retail financial services mis-selling scandal in UK history. It was sold to millions of customers who did not need it, did not want it, and in many cases were not eligible to claim on it. It was frequently added to policies without customers' knowledge or informed consent. And it generated enormous profits for the institutions that sold it — at extraordinary cost to the customers who bought it.}
\paragraph{The total redress paid to affected customers \textbf{exceeded £50 billion} — a figure that dwarfs any other consumer financial services redress programme in UK history.}

\subsection{The Risk Assessment Failures}
\paragraph{\textbf{The incentive structure failure}. PPI was extraordinarily profitable — commission rates of 50-70\% of the premium were common, and the attachment of PPI to a credit product could double or more the financial value of that transaction to the selling institution. The incentive for front-line sales staff to sell PPI was overwhelming — and the financial consequences of not selling it, in a target-driven sales culture, were significant}
\paragraph{As discussed in Chapter~\ref{chap:People and Conduct Risk Assessment}, incentive structures that reward sales volume without adequate attention to the quality and appropriateness of what is sold are a direct and foreseeable cause of mis-selling. The PPI scandal was not primarily a failure of individual integrity — it was a failure of institutional design. A product architecture, a commission structure, and a sales culture that made mis-selling the path of least resistance for front-line staff was a conduct risk that was identifiable, assessable, and manageable — and that was not adequately identified, assessed, or managed.}
\paragraph{\textbf{The product design failure}. PPI policies were frequently designed with exclusions that made them worthless for the customers most likely to need them — the self-employed, those with pre-existing medical conditions, and those in temporary employment were commonly excluded from coverage that they had been sold on the basis of needing. A genuine product risk assessment — asking honestly whether the product delivered value to the customers for whom it was designed — would have identified this fundamental inadequacy.}
\paragraph{The Consumer Duty's requirement for genuine fair value assessment — asking whether the benefits of a product are reasonable relative to its cost for the customers who buy it — is in many respects a direct regulatory response to the PPI scandal. If that standard had been applied to PPI, the product could not have been sold as it was.}
\paragraph{\textbf{The complaints handling failure}. When customers complained about PPI — which they did, in large numbers — those complaints were systematically rejected by the selling institutions. Complaint handling processes that were designed to reject rather than to assess, that were measured on rejection rates rather than on outcomes for customers, and that were operated by staff with the same incentive misalignment as the original sales process, failed to provide the corrective feedback that a genuine risk monitoring system would have delivered.}
\paragraph{The pattern of complaints was a monitoring signal of enormous significance — a clear indicator that something was seriously wrong with the product and its sale. The failure to treat that signal as the risk intelligence it represented — the failure to ask \textit{"why are so many customers complaining, and what does that tell us about our product and our sales practices?"} — is a conduct risk monitoring failure of the first order}
\paragraph{\textbf{The regulatory intelligence failure}. The FSA — the FCA's predecessor — had identified concerns about PPI mis-selling as early as 2005 and had begun publishing guidance and undertaking supervisory work in this area. That regulatory intelligence — the signal that the regulator had identified the product as a conduct risk — should have been a powerful trigger for internal risk assessment and remediation. For many institutions, it was not.}

\subsection{The Lessons}
\paragraph{\textbf{Product design is a conduct risk assessment exercise}. Before any product is sold, a genuine assessment of whether it delivers fair value to the customers for whom it is intended — whether its design, its pricing, and its exclusions are consistent with the interests of those customers — is a fundamental conduct risk obligation. Products that are primarily designed to generate revenue for the seller rather than value for the customer are conduct risks, and should be assessed as such.}
\paragraph{\textbf{Incentive misalignment is a foreseeable and manageable conduct risk}. The connection between high-pressure, volume-driven sales incentives and mis-selling is not mysterious — it is predictable. An organisation that designs an incentive structure that rewards volume without adequate quality controls, and that fails to assess the conduct risk that structure creates, has failed in its most basic conduct risk management obligation.}
\paragraph{\textbf{Complaints are risk intelligence}. A pattern of customer complaints is one of the most direct and most valuable signals of conduct risk available to an organisation. Complaint handling processes that are designed to reject rather than to learn are a governance failure — they destroy the most important feedback mechanism the organisation has for identifying conduct problems before they become crises.}
\paragraph{\textbf{Regulatory intelligence is a risk trigger}. When the regulator signals concern about a practice or product — through guidance, through supervisory work, through enforcement action at peer firms — that signal is a direct trigger for internal risk assessment. Organisations that treat regulatory intelligence as an external compliance requirement rather than as risk information relevant to their own assessment are missing one of the most valuable inputs available.}

%---------------------------- SECTION ----------------------------
\section{Case 5 — Boeing 737 MAX: The Failure of Safety Culture Under Commercial Pressure}

\subsection{What Happened}
\paragraph{Between October 2018 and March 2019, two Boeing 737 MAX aircraft crashed — Lion Air Flight 610 and Ethiopian Airlines Flight 302 — killing 346 people. Both crashes were caused by the same system: the Maneuvering Characteristics Augmentation System (MCAS), an automated flight control feature that repeatedly pushed the aircraft's nose down in response to faulty sensor data, overwhelming the pilots' attempts to recover.}
\paragraph{The subsequent investigations — by the US Congress, the Department of Justice, the FAA, and multiple other bodies — revealed a systemic failure of safety risk assessment at Boeing and in the regulatory approval process, driven by intense commercial pressure to bring the MAX to market quickly and cheaply.}

\subsection{The Risk Assessment Failures}
\paragraph{\textbf{The commercial pressure failure}. Boeing was under intense pressure to respond to the competitive threat posed by the Airbus A320neo — a fuel-efficient narrow-body aircraft that was attracting significant airline orders. Rather than designing a new aircraft, Boeing chose to reengine the existing 737 platform — but the larger, more fuel-efficient engines changed the aircraft's flight characteristics in ways that required MCAS to compensate. The pressure to deliver the MAX as a 737 variant — rather than as a new type requiring more extensive pilot retraining and recertification — drove decisions about MCAS design and certification that prioritised schedule and cost over safety.}
\paragraph{This is the commercial pressure dynamic at its most consequential — the systematic distortion of risk assessment by the need to reach conclusions that support a commercially desirable outcome. As discussed in Chapter~\ref{chap:Communicating Risk Across the Organisation} in the context of groupthink, the pressure to confirm rather than challenge is most powerful when the stakes — commercial, reputational, competitive — are highest. It is precisely in those conditions that the independence and rigour of risk assessment most needs to be protected.}
\paragraph{\textbf{The MCAS design and disclosure failure}. MCAS was significantly more powerful than had been disclosed to airlines and pilots during the type certification process. Its authority — the degree to which it could deflect the horizontal stabiliser — had been increased during development, but the safety assessment had not been updated to reflect the increased authority. Pilots were not informed of MCAS's existence in the flight manual — in part because disclosing it would have triggered additional training requirements that would have undermined the commercial case for the MAX as a 737 variant.}
\paragraph{A safety risk assessment that does not reflect the actual design of the system being assessed — that is based on an earlier, less powerful version of a critical safety system — is not a genuine safety assessment. It is a document whose primary function is regulatory approval rather than safety assurance.}
\paragraph{\textbf{The regulatory capture failure}. The FAA's certification of the 737 MAX involved significant delegation of the certification work to Boeing employees acting as Authorised Representatives of the FAA — a delegation model that created structural conflicts of interest. Employees who were responsible for Boeing's commercial schedule were also responsible for conducting safety assessments on Boeing's behalf. The oversight function had, in effect, been captured by the organisation it was supposed to oversee.}
\paragraph{For compliance professionals, regulatory capture — the progressive alignment of regulatory oversight with the interests of the regulated rather than those of the public — is a risk that applies to internal oversight functions as much as to external regulators. A compliance function that becomes too aligned with commercial objectives, that avoids challenging uncomfortable conclusions, and that measures its success by its facilitation of business rather than its rigour of oversight, has undergone the same kind of capture.}
\paragraph{\textbf{The pilot training failure}. When MCAS malfunctioned, the pilots of both crashed aircraft were not adequately prepared to diagnose and respond to the situation. The decision to minimise pilot retraining requirements — for commercial reasons — meant that pilots encountered a novel and dangerous flight control failure without the knowledge or the simulator training to manage it. The risk of this training gap was identifiable — and was identified, privately, by some Boeing employees — but the commercial pressure to avoid the training trigger was not adequately challenged.}

\subsection{The Lessons}
\paragraph{\textbf{Commercial pressure must not distort safety risk assessment}. The most dangerous risk assessment environments are those where the conclusion of the assessment is implicitly predetermined — where the pressure to reach a commercially necessary conclusion is so powerful that the assessment process becomes, in substance, a documentation exercise rather than a genuine inquiry. Protecting the independence and the rigour of safety risk assessment from commercial pressure is a governance obligation, not just a cultural aspiration.}
\paragraph{\textbf{Safety assessments must reflect the actual design, not the approved design}. A risk assessment that is based on an earlier, less risky version of a product, system, or process — that is not updated to reflect subsequent design changes that increase risk — is not a genuine assessment. Change management processes must include a clear mechanism for triggering reassessment when changes are made that affect the risk profile.}
\paragraph{\textbf{Delegated oversight requires robust independence protections}. Where oversight functions are delegated to the organisation being overseen — whether in regulatory certification, in third-party assurance, or in internal audit — the independence of the oversight function from the commercial interests of the organisation must be actively protected. Structural conflicts of interest in oversight arrangements are themselves significant governance risks.}

%---------------------------- SECTION ----------------------------
\section{Case 6 — Deepwater Horizon: Multiple Barrier Failure}

\subsection{What Happened}
\paragraph{On 20 April 2010, the Deepwater Horizon drilling rig in the Gulf of Mexico suffered a blowout — an uncontrolled release of oil and gas from the well — resulting in an explosion and fire that killed eleven workers and injured seventeen others. The rig sank two days later, and the resulting oil spill — 4.9 million barrels over 87 days before the well was capped — was the largest accidental marine oil spill in history.}
\paragraph{The subsequent investigations identified a catastrophic failure of multiple safety barriers — the physical, procedural, and organisational systems that were supposed to prevent exactly this kind of event.}

\subsection{The Risk Assessment Failures}
\paragraph{\textbf{The Swiss Cheese alignment}. The Deepwater Horizon disaster is perhaps the clearest real-world illustration of Reason's Swiss Cheese model. Multiple independent safety barriers — the well's cement job, the negative pressure test, the blowout preventer, and the emergency disconnect system — all failed, or were misinterpreted, or were not activated, in a sequence that allowed the blowout to proceed to catastrophe. No single failure was sufficient to cause the disaster — but the alignment of failures across all the defensive layers created the pathway through which the hazard became harm.}
\paragraph{The lesson of the Swiss Cheese model is not that each barrier needs to be perfect — it is that the design, maintenance, and testing of the barrier system as a whole needs to be adequate to prevent the alignment of failures that causes catastrophic outcomes. An organisation that manages each barrier independently — without assessing the combined adequacy of the barrier system — is systematically blind to the alignment risk.}
\paragraph{\textbf{The negative pressure test misinterpretation}. In the hours before the blowout, a negative pressure test — designed to confirm the integrity of the well — produced anomalous results that indicated the well was not secure. Those results were misinterpreted by the drilling team as satisfactory — a misinterpretation that was contested by some individuals on the rig but that was not escalated or resolved before operations continued.}
\paragraph{The negative pressure test failure illustrates the point-of-work risk assessment failure — the failure to recognise that current conditions differ materially from those anticipated in the standing risk assessment, and to respond appropriately. The anomalous test results were a real-time signal of significant risk that was misread, in part because of time pressure to complete the well, in part because of overconfidence in the team's interpretation, and in part because the organisational culture did not adequately support the kind of persistent, uncomfortable challenge that the situation required.}
\paragraph{\textbf{The cost and time pressure}. The Deepwater Horizon well was running significantly behind schedule and over budget at the time of the disaster. The decisions that contributed most directly to the disaster — the reduced number of centralizers in the cement job, the decision not to run a cement bond log to verify the cement job, the decision to proceed despite the anomalous negative pressure test results — were all decisions that saved time and money at the cost of safety margin.}
\paragraph{The pattern — commercial pressure driving progressive reduction of safety margins without adequate risk assessment of the cumulative effect — is precisely the dynamic of normalisation of deviance described in the Challenger case. The individual decisions each appeared marginal — each one saved only a small amount of time or money and was assessed in isolation as acceptable. Their cumulative effect was not assessed.}
\paragraph{\textbf{The inadequate testing of the blowout preventer}. The blowout preventer — the last line of defence against a blowout — had not been adequately tested to confirm it would function in the conditions of the Deepwater Horizon well. A safety barrier that is not tested is a safety barrier whose reliability is unknown — and an unknown reliability is not the same as a confirmed adequate reliability.}

\subsection{The Lessons}
\paragraph{\textbf{Safety barrier systems must be assessed as systems}, not as individual components. The adequacy of a multi-layered safety barrier system depends not just on the adequacy of each individual barrier but on the combined adequacy of the system — including its resilience to the simultaneous failure of multiple barriers. Risk assessment needs to address the system, not just its components.}
\paragraph{\textbf{Anomalous signals during operations must be treated as risk alerts, not problems to be explained away}. The negative pressure test results were anomalous — they required a satisfactory explanation before operations could responsibly continue. The failure to require such an explanation — to treat the anomaly as something that needed to be resolved rather than reinterpreted — reflects a culture in which schedule pressure distorts the response to safety signals.}
\paragraph{\textbf{Cumulative margin erosion must be assessed in aggregate}. Individual decisions that each appear marginally acceptable may be collectively unacceptable when their combined effect on safety margin is assessed. Risk assessment processes need to track the cumulative effect of decisions that individually appear marginal — recognising that the accumulation of small compromises can produce a catastrophic reduction in overall safety margin.}

%---------------------------- SECTION ----------------------------
\section{Case 7 — The Weinstein Company: Governance Failure and Conduct Risk}

\subsection{What Happened}
\paragraph{The revelations about Harvey Weinstein — which became public in October 2017 through reporting in the New York Times and The New Yorker — triggered one of the most significant corporate governance and conduct risk failures of the decade. Weinstein had engaged in sexual harassment and assault of employees and others over decades — conduct that was known to, or suspected by, a significant number of people within and around the Weinstein Company, but that was systematically concealed, enabled, and in some cases actively covered up.}
\paragraph{The corporate failure that the Weinstein case illustrates is not primarily a legal or compliance failure — it is a governance and conduct risk failure of a fundamental character.}

\subsection{The Risk Assessment Failures}
\paragraph{\textbf{The board oversight failure}. The board of the Weinstein Company — despite evidence that significant conduct concerns existed — did not exercise adequate oversight of the conduct risks associated with its founder and chief executive. The power imbalance between a founder-CEO with complete control of the organisation and a board that was dependent on his creative and commercial success for the company's viability created conditions in which genuinely independent oversight was structurally difficult.}
\paragraph{Board oversight of conduct risk — particularly where the conduct risk is most acute in relation to the most powerful individual in the organisation — is one of the most demanding governance challenges. It requires board members who are genuinely independent, genuinely willing to challenge, and genuinely empowered to act — and a governance structure that provides them with the information they need to exercise that oversight.}
\paragraph{\textbf{The settlement and concealment culture}. Multiple complaints and settlements — involving non-disclosure agreements — had been made over the years before the public revelations. Each settlement resolved an individual complaint without addressing the underlying conduct — and the non-disclosure agreements actively suppressed the accumulation of evidence that might have revealed the systemic nature of the problem.}
\paragraph{The use of settlement and non-disclosure as a primary response to conduct complaints — rather than as a last resort after genuine investigation and remediation — is a conduct risk governance failure. It treats conduct risk as a liability management problem rather than as a genuine risk requiring assessment and management. And it destroys the monitoring intelligence that a genuine conduct risk framework would use to detect systemic problems.}
\paragraph{\textbf{The cultural enablement}. The conduct was possible — and persistent — because the organisational culture enabled it. Power dynamics that silenced those who experienced or witnessed the conduct, a culture in which the founder's behaviour was treated as beyond challenge, and an environment in which those who raised concerns faced professional consequences — all created the conditions in which serious conduct risk accumulated without effective check.}
\paragraph{As discussed in Chapter~\ref{chap:People and Conduct Risk Assessment}, conduct risk is fundamentally a cultural and structural phenomenon — it arises from the interaction of individual behaviour, power dynamics, incentive structures, and the degree to which organisational culture genuinely holds people to account regardless of their commercial value. Risk assessment that addresses only the formal governance structures without engaging honestly with the cultural conditions is systematically blind to the most significant conduct risks.}

\subsection{The Lessons}
\paragraph{\textbf{Conduct risk assessment must address power dynamics explicitly}. The most serious conduct risks are often concentrated in precisely the areas where power dynamics make challenge most difficult — around founders, around high performers, around individuals whose commercial value has been allowed to translate into immunity from accountability. Conduct risk assessment needs to be explicitly attentive to these concentrations of power and the specific governance challenges they create.}
\paragraph{\textbf{Settlement is not remediation}. Managing conduct complaints through settlement and non-disclosure — without genuine investigation, genuine accountability, and genuine cultural change — does not reduce conduct risk. It suppresses the monitoring intelligence that would reveal the systemic nature of the problem whilst allowing the underlying conditions to persist.}
\paragraph{\textbf{Non-disclosure agreements can be risk governance tools or risk concealment tools. Their use as a primary response to conduct complaints — rather than as an occasional last resort — transforms them from legitimate legal instruments into mechanisms for suppressing the risk intelligence that governance bodies need.}}

%---------------------------- SECTION ----------------------------
\section{Case 8 — Barings Bank: The Limits of Oversight in Complex Operations}

\subsection{What Happened}
\paragraph{In February 1995, Barings Bank — the oldest merchant bank in the UK — collapsed after a single trader, Nick Leeson, accumulated losses of £827 million through unauthorised trading in derivatives on Japanese equity indices. Leeson had been able to accumulate these losses because he controlled both the front office trading desk and the back office settlement function in Barings' Singapore operation — a fundamental breach of the segregation of duties principle — and because the oversight mechanisms that should have detected his activities had failed comprehensively.}

\subsection{The Risk Assessment Failures}
\paragraph{\textbf{The segregation of duties failure}. The most fundamental control failure was the absence of segregation between front office and back office functions. Leeson's ability to control both his own trading and the settlement and reporting of that trading gave him the ability to conceal his losses through a fictitious account — and to do so for more than two years without detection. This is precisely the structural vulnerability that segregation of duties is designed to prevent — and its absence was a basic operational risk assessment failure.}
\paragraph{\textbf{The oversight failure}. Leeson was generating exceptional profits — or appeared to be — and those exceptional profits suppressed the scrutiny that should have been applied to his operations. The apparent success of the Singapore operation created a cultural dynamic in which questioning Leeson's activities was unwelcome — he was a star performer, and challenging a star performer requires both the institutional will and the governance structure to support that challenge.}
\paragraph{This dynamic — in which commercial success suppresses risk oversight — is a recurring theme across conduct and fraud cases. Organisations that apply less scrutiny to high-performing areas than to lower-performing ones create a systematic blind spot precisely in the areas where the incentive for misconduct may be highest.}
\paragraph{\textbf{The risk assessment of a novel business}. Barings' Singapore operation was conducting a business — proprietary trading in derivatives — that was relatively new, technically complex, and not well understood by the senior management responsible for overseeing it. Risk assessment of activities that are not well understood is inherently challenging — and the temptation to substitute trust in apparently successful individuals for genuine understanding of the risks they are taking is strong.}
\paragraph{\textbf{The management information failure}. The management information flowing from Singapore to Barings' London head office was inadequate — it did not provide the visibility of Leeson's actual risk positions that would have enabled early detection of the problem. Risk assessment requires adequate information — and organisations that do not invest in the management information systems needed to maintain visibility of significant risk positions are flying blind.}

\subsection{The Lessons}
\paragraph{\textbf{Segregation of duties is not merely a best-practice recommendation}. It is one of the most fundamental operational risk controls available — and its absence in any context involving significant financial or asset risk is a material governance failure. No single individual should have control of both the initiation and the recording or settlement of significant transactions.}
\paragraph{\textbf{Exceptional performance warrants more scrutiny, not less}. Returns that are materially better than peers, operations that appear consistently profitable without explanation of the source of that profitability, and individuals who resist oversight whilst consistently delivering results — all should attract more scrutiny, not less. The appearance of success is not evidence of the absence of risk.}
\paragraph{\textbf{Adequate management information is a risk assessment prerequisite}. Organisations cannot assess risks they cannot see. Investment in management information systems that provide genuine visibility of significant risk positions — in financial operations, in compliance activities, and in any other context where risk accumulates — is a prerequisite for effective risk assessment, not an optional upgrade.}

%---------------------------- SECTION ----------------------------
\section{The Recurring Patterns — What the Cases Tell Us}
\paragraph{Across the eight cases examined in this chapter — and across the many more that could have been chosen — several recurring patterns of risk assessment failure emerge with striking consistency.}

\subsection{Pattern 1 — The Incentive Distortion}
\paragraph{In the 2008 crisis, in PPI, in the Boeing 737 MAX, and in Barings Bank, incentive structures that rewarded risk-taking, commercial performance, or specific behaviours without adequate accountability for risk outcomes distorted the risk assessment process — creating systematic pressure towards optimistic assessments and away from honest acknowledgement of risk.}
\paragraph{The lesson is simple but demanding: \textbf{incentive structures must be assessed as risk factors}. An organisation whose incentive structures are misaligned with its risk management objectives has a structural conduct and operational risk that no amount of policy or procedure can fully mitigate.}

\subsection{Pattern 2 — The Commercial Pressure Override}
\paragraph{In Challenger, Boeing, and Deepwater Horizon, commercial pressure — schedule pressure, cost pressure, competitive pressure — distorted the risk assessment and decision-making process in ways that produced catastrophic outcomes. The pattern is consistent: the pressure to reach a commercially necessary conclusion undermines the independence of the risk assessment, and the safety margins that are progressively eroded by commercially driven decisions accumulate into a systemic vulnerability that is not assessed in aggregate.}
\paragraph{The lesson: \textbf{commercial pressure is the most reliable predictor of risk assessment distortion}. Governance structures that explicitly protect the independence of risk assessment from commercial pressure — that require challenges to be resolved rather than overridden, that make the override of technical safety concerns a documented and individually accountable decision — are the most effective defence.}

\subsection{Pattern 3 — The Normalisation of Deviance}
\paragraph{In Challenger, in Deepwater Horizon, and in the pattern of conduct tolerance in the Weinstein case, the progressive normalisation of deviation from required standards — each individual deviation appearing manageable in isolation — accumulated into systemic risk that was not visible from any individual instance.}
\paragraph{The lesson: \textbf{risk assessment must monitor for cumulative deviation, not just individual incidents}. A culture that tolerates small deviations from standards because they have not yet produced adverse outcomes is a culture in which normalisation of deviance is underway — and in which the risk register's failure to capture the cumulative pattern is a significant governance gap.}

\subsection{Pattern 4 — The Communication and Escalation Failure}
\paragraph{In Challenger, in Grenfell, and in Barings, critical risk information existed — was known to individuals who understood its significance — but failed to reach the decision-makers who had the authority and the responsibility to act on it. The information existed; the communication and escalation structures failed.}
\paragraph{The lesson: \textbf{risk assessment quality is determined by information flow as much as by analytical technique}. Governance structures that protect the ability of those with relevant risk knowledge to communicate it — that create safe channels for escalation, that genuinely act on concerns raised, and that hold decision-makers accountable for the decisions they make with the information available to them — are as important as the risk assessment methodology itself.}

\subsection{Pattern 5 — The False Assurance of Documentation}
\paragraph{Across virtually every case examined — from the risk models of the financial crisis, to the test certificates of Grenfell, to the safety assessments of Boeing — the existence of documentation, of certification, of apparently adequate formal processes created a false assurance that genuine risk management was occurring. The documentation existed; the substance did not.}
\paragraph{The lesson — stated repeatedly throughout this book and never more powerfully illustrated than in these cases — is the most important one: \textbf{the test of a risk assessment framework is not its documentation. It is what actually happens}.}

%---------------------------- SECTION ----------------------------
\newpage
\section{Chapter Summary}
In this chapter we learned about the following key points:

\begin{table}[H]
  \centering
  \renewcommand{\arraystretch}{1.5}
  \begin{tabularx}{\textwidth}{ | >{\raggedright\arraybackslash}p{0.20\textwidth} 
								| >{\raggedright\arraybackslash}p{0.44\textwidth} 
                                | >{\raggedright\arraybackslash}X | }
    \hline
    \textbf{Case} & \textbf{Primary Failure Type} & \textbf{Key Lesson} \\
    \hline
    \textbf{2008 Financial Crisis} & Model uncertainty; incentive distortion; complexity opacity & Quantitative confidence ≠ analytical reliability; incentives distort risk assessment\\
    \hline
    \textbf{Challenger Disaster} & Normalisation of deviance; communication failure; burden of proof reversal & Deviance without harm is not safe deviance; technical voice must be protected\\
    \hline
    \textbf{Grenfell Tower} & Regulatory gap exploitation; fragmented accountability; resident voice ignored & Compliance ≠ safety; those closest to risk have the most valuable intelligence\\
    \hline
    \textbf{PPI Mis-Selling} & Incentive misalignment; product design failure; complaints intelligence ignored & Incentive structures are conduct risks; complaints are risk intelligence\\
    \hline
    \textbf{Boeing 737 MAX} & Commercial pressure override; design change without reassessment; regulatory capture & Commercial pressure is the most reliable predictor of risk assessment distortion\\
    \hline
    \textbf{Deepwater Horizon} & Swiss Cheese alignment; cumulative margin erosion; anomalous signal misinterpretation & Barrier systems must be assessed as systems; cumulative margin erosion is a risk\\
    \hline
    \textbf{Weinstein Company} & Power dynamics; settlement as concealment; cultural enablement & Conduct risk is concentrated where power makes challenge hardest\\
    \hline
    \textbf{Barings Bank} & Segregation failure; success-suppressed scrutiny; management information gap & Exceptional performance warrants more scrutiny, not less\\
    \hline
  \end{tabularx}
  \caption{Failures - The Examined Cases}
\end{table}

\begin{table}[H]
  \centering
  \renewcommand{\arraystretch}{1.5}
  \begin{tabularx}{\textwidth}{ | >{\raggedright\arraybackslash}p{0.3\textwidth} 
								| >{\raggedright\arraybackslash}p{0.35\textwidth} 
                                | >{\raggedright\arraybackslash}X | }
    \hline
    \textbf{Pattern} & \textbf{Description} & \textbf{Governance Response} \\
    \hline
    \textbf{Incentive distortion} & Reward structures drive optimistic risk assessment & Assess incentive structures as conduct risks\\
    \hline
    \textbf{Commercial pressure override} & Schedule/cost pressure undermines assessment independence & Protect assessment independence; document overrides with individual accountability\\
    \hline
    \textbf{Normalisation of deviance} & Cumulative deviation not assessed in aggregate & Monitor deviation patterns, not just adverse outcomes\\
    \hline
    \textbf{Communication and escalation failure} & Risk knowledge fails to reach decision-makers & Protect escalation channels; act on concerns raised\\
    \hline
    \textbf{False assurance of documentation} & Process existence substitutes for substance & Test what actually happens, not what documentation says\\
    \hline
  \end{tabularx}
  \caption{Failures - The 5 Recurring Patterns}
\end{table}

\paragraph{Key takeaways:}
\begin{itemize}
  \item{Real cases demonstrate that the most serious risk failures are rarely caused by ignorance of risk — they are caused by the systematic distortion of risk assessment by incentives, commercial pressure, power dynamics, and cognitive bias.}
  \item{The organisations most damaged by risk failures frequently had risk assessment frameworks that were, on paper, adequate — the failures were in the culture, the governance, and the incentive structures that determined whether those frameworks genuinely operated.}
  \item{Normalisation of deviance — the progressive acceptance of deviation without adverse consequence as evidence of safety — is one of the most powerful and most insidious dynamics in organisational risk failure.}
  \item{Exceptional performance should attract more scrutiny, not less — it is one of the most reliable indicators of either exceptional capability or concealed risk.}
  \item{The most valuable risk intelligence often resides with those closest to the risk — employees, customers, residents — and the failure to engage with and act on that intelligence is both a governance failure and, as Grenfell demonstrated, a moral one.}
  \item{The recurring lesson of real cases — stated at the beginning of this book and confirmed by every case in this chapter — is that the test of a risk assessment framework is not its documentation but what actually happens.}
\end{itemize}

\barquote{In Chapter~\ref{chap:Designing a Risk Assessment Programme for Your Organisation}, we bring everything together in the most practically oriented chapter of the book — Designing a Risk Assessment Programme for Your Organisation. Drawing on all the frameworks, methodologies, and lessons of the preceding chapters, we provide a comprehensive, actionable guide to building a risk assessment programme that is genuinely effective, proportionate, and sustainable.}